% D1-grav-waves  D16+D24+D26
% Auto-generated from markdown source
% origin: ai

\documentclass[12pt, a4paper]{article}
\usepackage[utf8]{inputenc}
\usepackage[T1]{fontenc}
\usepackage{amsmath, amssymb, amsfonts, amsthm}
\usepackage{textcomp}
\usepackage{geometry}
\usepackage{hyperref}
\usepackage{booktabs}
\usepackage{graphicx}
\usepackage{fancyhdr}
\usepackage{tcolorbox}
\usepackage{enumitem}

\geometry{margin=1in}
\pagestyle{fancy}
\fancyhf{}
\rhead{D1-grav-waves  D16+D24+D26}
\lhead{\thepage}

\newtheorem{definition}{Definition}[section]
\newtheorem{theorem}{Theorem}[section]
\newtheorem{corollary}{Corollary}[section]
\newtheorem{proposition}{Proposition}[section]
\newtheorem{lemma}{Lemma}[section]
\newtheorem{conjecture}{Conjecture}[section]

\newcommand{\calI}{\mathcal{I}}
\newcommand{\calP}{\mathcal{P}}
\newcommand{\calE}{\mathcal{E}}
\newcommand{\calS}{\mathcal{S}}
\newcommand{\calH}{\mathcal{H}}
\newcommand{\calO}{\mathcal{O}}
\newcommand{\calF}{\mathcal{F}}
\newcommand{\calU}{\mathcal{U}}
\newcommand{\calV}{\mathcal{V}}
\newcommand{\calA}{\mathcal{A}}
\newcommand{\calG}{\mathcal{G}}
\newcommand{\calB}{\mathcal{B}}
\newcommand{\calC}{\mathcal{C}}
\newcommand{\calL}{\mathcal{L}}
\newcommand{\calM}{\mathcal{M}}
\newcommand{\calX}{\mathcal{X}}
\newcommand{\calY}{\mathcal{Y}}
\newcommand{\calW}{\mathcal{W}}
\newcommand{\calT}{\mathcal{T}}
\newcommand{\calZ}{\mathcal{Z}}
\newcommand{\calR}{\mathcal{R}}
\newcommand{\calQ}{\mathcal{Q}}
\newcommand{\calD}{\mathcal{D}}
\newcommand{\calN}{\mathcal{N}}
\newcommand{\calK}{\mathcal{K}}

\newcommand{\Tr}{\text{Tr}}
\newcommand{\E}{\mathbb{E}}
\newcommand{\diff}{\mathrm{d}}
\newcommand{\KL}{D_{\mathrm{KL}}}
\newcommand{\argmax}{\arg\max}
\newcommand{\argmin}{\arg\min}
\newcommand{\hbarF}{\hbar}

\newtcolorbox{thmbox}[1]{
  colback=blue!5!white, colframe=blue!75!black, fonttitle=\bfseries,
  title={#1}, breakable
}
\newtcolorbox{defbox}[1]{
  colback=green!5!white, colframe=green!75!black, fonttitle=\bfseries,
  title={#1}, breakable
}

\begin{document}

\title{D1-grav-waves  D16+D24+D26}
\author{RALPH Loop}
\date{}
\maketitle

\begin{abstract}
This paper presents D1-grav-waves  D16+D24+D26, formalizing the underlying principles using information-theoretic methods.

\textbf{Keywords}: D1-grav-waves  D16+D24+D26, information theory, formalization.
\end{abstract}

\tableofcontents
\section{D1-grav-waves  D16+D24+D26}



\subsection{0.}


\textbf{D1-grav-waves } Extended Functorial QFT **** Hellings-Downs / postulated

\textbf{D16+D24+D26 } Fisher ****

****



\subsection{1.}


\subsubsection{1.1 D1-grav-waves}


|  | D1-grav-waves |  |
|------|---------------|------|
| 1 |  | Extended Functorial AQFT \(F: \mathbf{Loc}_{M}^{\mathrm{ext}} \to \mathbf{C^*Alg}\) |
| 2 |  |  \(w_i\)  \(R_i(t)\) |
| 3 |  |  cobordism  |
| 4 |  | HD  =  3.1  |
| 5 |  |  =  |

\subsubsection{1.2 D16+D24+D26}


|  | D16+D24+D26 |  |
|------|-------------|------|
| 1 |  | Fisher  \((\Sigma, g_F)\) |
| 2 |  |  \(\sigma\)  \(p(x|\sigma)\) |
| 3 |  |  \(\mathcal{C}\)  \(q_{\mathcal{C}}(x|\sigma)\) |
| 4 |  |  = Fisher  = KL  1 |
| 5 |  |  \(I(\mathcal{C}; T)\) =  |

\subsubsection{1.3}


D1  "ALGEBRAIC STATE DISTANCE  SOURCE DISCRIMINATION"  D16+D24  "FISHER METRIC DISTANCE  SEMANTIC CLASSIFICATION"

****

\[\begin{array}{c|c}
\text{D1-grav-waves} & \text{D16+D24+D26} \\
\hline
\text{Algebraic state } \omega & \text{Parameter } \sigma \in \Sigma \\
\text{State distance } d_{TV}(\omega_1, \omega_2) & \text{Semantic distance } d_{g_F}(\sigma_1, \sigma_2) \\
\text{Source class} & \text{Classification } \mathcal{C} \\
\text{Total variation distance} & \text{KL divergence / mutual information} \\
\text{Cobordism structure} & \text{Fisher metric structure} \\
\text{HD correlation function} & \text{Semantic geodesic} \\
\text{NANOGrav data} & \text{NLP concept embeddings}
\end{array}\]

\subsection{2.}


\subsubsection{2.1 Distance Derivation Principle}


****
\(\mathcal{S}\) /****fundamental structure\(\mathcal{O}\) \(d\)  \(\mathcal{O}\)

**** \(\Phi: \mathcal{O} \to \mathcal{S}\)  \(\Phi\)  \(\Phi\)  \(d\)  \(\mathcal{S}\)

\textbf{ D1 }
\begin{itemize}
\item \(\mathcal{S}\) = Extended Functorial AQFT
\item \(\mathcal{O}\) = PTA
\item \(\Phi\) =  \(w_i \mapsto \mathcal{A}(w_i)\)
\item \(d\) = HD
\end{itemize}


\textbf{ D16+D24 }
\begin{itemize}
\item \(\mathcal{S}\) = Fisher
\item \(\mathcal{O}\) =
\item \(\Phi\) =  \(p(x|\sigma)\)
\item \(d\) =  KL
\end{itemize}


\subsection{3.}


\subsubsection{3.1  D26  D1-grav-waves}


D1-grav-waves  \(d_{TV}(\omega_{\text{SMBHB}}, \omega_{\text{cosmo}})\)

**** NANOGrav

\textbf{D26 } Fisher \(\min_i \text{Re}(\lambda_i) \to 0\)""

\subsubsection{3.2  D1  D16+D24}


D1  D16+D24 ****

D16+D24  KL
\begin{itemize}
\item "" \(\mathbf{Sem}\)
\item "" \(\mathbf{Stat}\)
\item \(G: \mathbf{Sem} \to \mathbf{Stat}\)
\item = \(G\)  \(\mathbf{Stat}\)
\end{itemize}


\subsubsection{3.3}


\textbf{*}*** \(\mathbf{DistCat}\) D1-grav-waves  D16+D24+D26

\textbf{ \(\mathbf{DistCat}\)}
\(\mathbf{DistCat}\)  2-
\begin{itemize}
\item 0-
\item 1-
\item 2-
\item = 2-
\end{itemize}


\textbf{D1-grav-waves }
\(\mathbf{DistCat}\)  0- = Extended Functorial AQFT
\(\mathbf{DistCat}\)  1- =  \(F: \mathbf{Loc} \to \mathbf{C^*Alg}\)
\(\mathbf{DistCat}\)  2- =  \(w_i, w_j \mapsto \langle R_i R_j \rangle\)
= HD

\textbf{D16+D24 }
\(\mathbf{DistCat}\)  0- = Fisher
\(\mathbf{DistCat}\)  1- =  \(p(x|\sigma)\)
\(\mathbf{DistCat}\)  2- =  \(q_{\mathcal{C}}(x|\sigma)\)
=

\subsection{4.}


1. ****D1-grav-waves  D16+D24+D26 ""
2. \textbf{ 1}D26  D1-grav-waves
3. \textbf{ 2}D1  D16+D24
4. **** HD

\subsection{5.}


\begin{itemize}
\item [ ]  \(\mathbf{DistCat}\)
\item [ ]  \(\mathbf{DistCat}\)  D1  HD  D16+D24
\item [ ]  NLP
\end{itemize}


\begin{thebibliography}{99}
\bibitem{ref1} Shannon, C.E. (1948). A Mathematical Theory of Communication.
\bibitem{ref2} Cover, T.M. (2006). Elements of Information Theory.
\end{thebibliography}

\end{document}